
\bta{2005}

\begin{enumerate}
	%\renewcommand{\labelenumi}{\arabic{enumi}.}
	% A(\Alph) a(\alph) I(\Roman) i(\roman) 1(\arabic)
	%设定全局标号series=example	%引用全局变量resume=example
	%[topsep=-0.3em,parsep=-0.3em,itemsep=-0.3em,partopsep=-0.3em]
	%可使用leftmargin调整列表环境左边的空白长度 [leftmargin=0em]
	\item
（$20'$）1800个电子经$1000V$电势差加速后从$x=-\infty$处射向势阶$V(x)=\begin{cases}V_0,&x<0\\0,&x>0\end{cases}$，其中$V_0=750V$。试问在$x=\infty$处能观察到多少个电子？如果势阶翻转一下，即电子射向势阶$V(x)=\begin{cases}0,&x<0\\V_0,&x>0\end{cases}$，则结果如何？


\banswer{
	电子经 $1000\,\mathrm{V}$ 加速后的入射动能为 $K_1=1000\,\mathrm{eV}$。阶梯两侧波数分别为 $k_i=\sqrt{2mK_i}/\hbar$，透射系数为
	\[
		T=\frac{4k_1k_2}{(k_1+k_2)^2}
		=\frac{4\sqrt{K_1K_2}}{(\sqrt{K_1}+\sqrt{K_2})^2}.
	\]
	对题中下降势阶，$K_2=K_1+750\,\mathrm{eV}=1750\,\mathrm{eV}$，故
	\[
		T=\frac{4\sqrt{1000\times1750}}{(\sqrt{1000}+\sqrt{1750})^2}
		\simeq0.98068.
	\]
	在 $x=+\infty$ 处平均可观察到 $1800T\simeq1765$ 个电子。

	势阶翻转后 $K_2=1000-750=250\,\mathrm{eV}$，于是
	\[
		T=\frac{4\sqrt{1000\times250}}{(\sqrt{1000}+\sqrt{250})^2}
		=\frac89,
	\]
	平均可观察到 $1800T=1600$ 个电子。
}


\item 
（$20'$）质量为$m$，电荷为$q$的粒子在三维各向同性谐振子势$V(r)=\frac{1}{2}m\omega^2r^2$中运动，同时受到一个沿$x$方向的均匀常电场$\vec{E}=E_0\vec{i}$作用。求粒子的能量本征值和第一激发态的简并度。此时轨道角动量是否守恒？如回答是，则请写出此守恒力学量的表达式。


\banswer{
	粒子的势能为
	\begin{align*}
		V&=\frac12m\omega^2(x^2+y^2+z^2)-qE_0x\\
		&=\frac12m\omega^2\left[(x-x_0)^2+y^2+z^2\right]
		-\frac{q^2E_0^2}{2m\omega^2},
		\qquad x_0=\frac{qE_0}{m\omega^2}.
	\end{align*}
	因此它只是以 $(x_0,0,0)$ 为新中心的各向同性谐振子，能量本征值为
	\[
		\boxed{E_{n_xn_yn_z}=\hbar\omega
		\left(n_x+n_y+n_z+\frac32\right)
		-\frac{q^2E_0^2}{2m\omega^2}}.
	\]
	第一激发态满足 $n_x+n_y+n_z=1$，空间简并度仍为 $3$。

	关于原点的轨道角动量 $\vec L=\vec r\times\vec p$ 不守恒，因为均匀电场产生力矩。但关于新势阱中心的
	\[
		\boxed{\vec L'=(\vec r-x_0\vec i)\times\vec p}
	\]
	满足 $[H,\vec L']=0$，是守恒力学量。
}


\item 
（$40'$）一个质量为$m$的粒子在下面的无限深方势阱中运动。$V(x)=\begin{cases}\infty,&x<0,x>a\\ 0,&a>x>0\end{cases}$开始时$(t=0)$，系统处于状态，$\psi(x)=A\sin\frac{\pi x}{2a}\cos^3\frac{\pi x}{2a}$，其中$A$为常数。请求出$t$时刻系统：
\begin{enumerate}
	%\renewcommand{\labelenumi}{\arabic{enumi}.}
	% A(\Alph) a(\alph) I(\Roman) i(\roman) 1(\arabic)
	%设定全局标号series=example	%引用全局变量resume=example
	%[topsep=-0.3em,parsep=-0.3em,itemsep=-0.3em,partopsep=-0.3em]
	%可使用leftmargin调整列表环境左边的空白长度 [leftmargin=0em]
	\item
处于基态的几率；
\item 
能量平均值；
\item 
动量平均值；
\item 
动量的均方差根（不确定度）。
	
\end{enumerate}


\banswer{
	利用
	\[
		\sin y\cos^3y=\frac14\sin2y+\frac18\sin4y,
		\qquad y=\frac{\pi x}{2a},
	\]
	并取 $\phi_n(x)=\sqrt{2/a}\sin(n\pi x/a)$，归一化后
	\[
		\psi(x,0)=\frac{2}{\sqrt5}\phi_1(x)+\frac{1}{\sqrt5}\phi_2(x),
		\qquad A=8\sqrt{\frac{2}{5a}}.
	\]
	因此
	\[
		\psi(x,t)=\frac{2}{\sqrt5}e^{-iE_1t/\hbar}\phi_1(x)
		+\frac{1}{\sqrt5}e^{-iE_2t/\hbar}\phi_2(x),
		\qquad E_n=\frac{n^2\pi^2\hbar^2}{2ma^2}.
	\]
	\begin{enumerate}
		\item 基态概率与时间无关：$P_1=4/5$。
		\item
		\[
			\langle H\rangle=\frac45E_1+\frac15E_2
			=\frac85E_1=\frac{4\pi^2\hbar^2}{5ma^2}.
		\]
		\item 需要保留两个能量分量之间的相干项。先计算
		\[
		 p_{12}=-\frac{4i\pi\hbar}{a^2}\int_0^a
		 \sin\frac{\pi x}{a}\cos\frac{2\pi x}{a}dx
		 =\frac{8i\hbar}{3a},\qquad p_{21}=p_{12}^*,\quad p_{11}=p_{22}=0.
		\]
		因此 $\langle p\rangle=(4/5)\operatorname{Re}(p_{12}e^{-i\Omega t})$。令 $\Omega=(E_2-E_1)/\hbar=3\pi^2\hbar/(2ma^2)$，可得
		\[
			\langle p\rangle_t=\frac{32\hbar}{15a}\sin(\Omega t).
		\]
		\item 因 $\langle p^2\rangle=(8/5)\pi^2\hbar^2/a^2$，
		\[
			\boxed{\Delta p(t)=\frac{\hbar}{a}
			\sqrt{\frac{8\pi^2}{5}-\frac{1024}{225}\sin^2(\Omega t)}}.
		\]
	\end{enumerate}
}


\newpage


\item 
（$30'$）两个具有相同质量$m$和频率$\omega$的谐振子，哈密顿量为
$$H^0=\frac{1}{2m}(p_1^2+p_2^2)+\frac{1}{2}m\omega^2((x_1-a)^2+(x_2+a)^2)$$
（$\pm a$为两个谐振子的平衡位置），受到微扰作用$H'=\lambda m\omega^2(x_1-x_2)^2,|\lambda|\ll1$，试求该体系的能级。


\banswer{
	令 $q_1=x_1-a$、$q_2=x_2+a$，再引入简正坐标
	\[
		Q=\frac{q_1+q_2}{\sqrt2},\qquad q=\frac{q_1-q_2}{\sqrt2}.
	\]
	对称模 $Q$ 的频率仍为 $\omega$；反对称模的势能为
	\[
		\frac12m\omega^2q^2+\lambda m\omega^2(\sqrt2q+2a)^2.
	\]
	把反对称模势能完整配方为
	\[
	 \frac12m\omega^2(1+4\lambda)
	 \left(q+\frac{4\sqrt2\lambda a}{1+4\lambda}\right)^2
	 +\frac{4\lambda m\omega^2a^2}{1+4\lambda}.
	\]
	正交坐标变换同时保持动能为两个独立的 $P^2/(2m)$，所以两模频率分别为
	$\omega$ 和 $\Omega=\omega\sqrt{1+4\lambda}$（要求 $1+4\lambda>0$）。
	故能级的精确形式为
	\[
		E_{n_+,n_-}=\hbar\omega\left(n_++\frac12\right)
		+\hbar\omega\sqrt{1+4\lambda}\left(n_-+\frac12\right)
		+\frac{4\lambda}{1+4\lambda}m\omega^2a^2.
	\]
	保留到 $\lambda$ 的一级，
	\[
		\boxed{E_{n_+,n_-}=\hbar\omega(n_++n_-+1)
		+2\lambda\hbar\omega\left(n_-+\frac12\right)
		+4\lambda m\omega^2a^2+O(\lambda^2)}.
	\]
	微扰使同一 $n_++n_-$ 多重态按反对称模量子数 $n_-$ 分裂。
}


\item 
（$30'$）已知氢原子基态波函数为：$\psi_{100}=\frac{1}{(\pi a_0^3)^{\frac{1}{2}}}e^{-\frac{r}{a_0}}$，试对坐标$x$及动量$p_x$，求：
$\Delta x=\sqrt{\langle x^2\rangle-\langle x^2\rangle},\Delta p=\sqrt{\langle p_x^2\rangle-\langle p_x\rangle^2}$，由此验证不确定关系。

\banswer{
	题中 $\Delta x$ 的定义应为 $\sqrt{\langle x^2\rangle-\langle x\rangle^2}$。基态具有球对称性，故
	\[
		\langle x\rangle=\langle p_x\rangle=0,\qquad
		\langle x^2\rangle=\frac13\langle r^2\rangle.
	\]
	将给定波函数和球坐标体积元代入，
	\[
	 \langle r^2\rangle=\frac4{a_0^3}\int_0^\infty r^4e^{-2r/a_0}dr
	 =\frac4{a_0^3}\frac{4!}{(2/a_0)^5}=3a_0^2,
	\]
	所以
	\[
		\Delta x=a_0.
	\]
	氢原子基态满足维里定理 $\langle T\rangle=-E_1=\hbar^2/(2ma_0^2)$，从而
	\[
		\langle p^2\rangle=2m\langle T\rangle=\frac{\hbar^2}{a_0^2},
		\qquad
		\langle p_x^2\rangle=\frac13\langle p^2\rangle
		=\frac{\hbar^2}{3a_0^2}.
	\]
	因此
	\[
		\Delta p_x=\frac{\hbar}{\sqrt3a_0},
		\qquad
		\boxed{\Delta x\,\Delta p_x=\frac{\hbar}{\sqrt3}>\frac\hbar2},
	\]
	验证了 Heisenberg 不确定关系。
}


\item 
考虑自旋 $\vec{s}$ 与角动量 $\hat{\vec{L}}$ 的耦合, 体系的哈密顿量为
\[
\hat{H}=-\frac{\hbar^{2}}{2 \mu} \nabla^{2}+V(r)+\lambda \hat{\vec{L}} \cdot \hat{\vec{S}}
\]
$\lambda$ 是耦合常数, 试证该体系的总角动量 $\hat{\vec{J}}=\hat{\vec{L}}+\hat{\overrightarrow{\mathrm{S}}}$ 守恒。\\
公式提示：在球坐标系内,
\[ 
\nabla^{2}=\frac{1}{r^{2}}\left(\frac{\partial}{\partial r}\left(r^{2} \frac{\partial}{\partial r}\right)-\frac{\vec{L}^{2}}{\hbar^{2}}\right), \quad \nabla f(r)=\frac{\vec{r}}{r} \frac{\partial}{\partial r} f(r), \quad \int_{0}^{\infty} t^{n} e^{-t} d t=n !
 \]

\banswer{
	球对称部分 $H_0=-\hbar^2\nabla^2/(2\mu)+V(r)$ 满足
	$[H_0,\vec L]=[H_0,\vec S]=0$。因此只需考察自旋--轨道项。对每个分量 $J_i=L_i+S_i$，
	\begin{align*}
		[\vec L\cdot\vec S,J_i]
		&=\sum_j\left([L_j,L_i]S_j+L_j[S_j,S_i]\right)\\
		&=i\hbar\sum_{jk}\left(\varepsilon_{jik}L_kS_j
		+\varepsilon_{jik}L_jS_k\right)=0,
	\end{align*}
	其中第二项交换哑指标 $j,k$ 后恰与第一项抵消。又因 $\lambda$ 为常数，
	\[
		[H,J_i]=[H_0,J_i]+\lambda[\vec L\cdot\vec S,J_i]=0.
	\]
	故 $d\vec J/dt=(i/\hbar)[H,\vec J]=0$，总角动量 $\vec J=\vec L+\vec S$ 守恒。
}


\end{enumerate}
